\ECchapter{05}{Energy Storage}{Why storing energy is key to the energy transition}{ECPurple}

\ECObjectives{
\begin{itemize}
\item Explain why electrical grids need energy storage.
\item Compare storage technologies using power, energy, duration and efficiency.
\item Select a suitable storage mix for a given engineering problem.
\end{itemize}}

\begin{multicols}{2}
\begin{engineeringnews}
As variable solar and wind generation expands, grid operators are deploying batteries,
pumped-storage plants and other forms of flexibility. The key challenge is not to find one
universal storage device, but to combine technologies operating over seconds, hours, days
and seasons.
\end{engineeringnews}

\begin{ecarticle}
Electricity networks must keep production and consumption balanced at every instant. A
sudden mismatch changes grid frequency and can damage equipment or trigger disconnections.
Storage systems absorb energy when production exceeds demand and return it when the system
needs additional power.

Engineers distinguish \textbf{power} from \textbf{energy}. Power, measured in watts,
describes how quickly a storage system can charge or discharge. Energy, measured in
watt-hours or joules, describes how long it can maintain that power. A flywheel may deliver
high power for a few seconds, whereas a reservoir or hydrogen cavern can store energy for
much longer periods.

Lithium-ion batteries provide rapid response and high round-trip efficiency. They are well
suited to frequency control, electric vehicles and daily solar shifting. Their limitations
include cost, ageing, fire protection and the supply of raw materials. Pumped-storage
hydroelectricity stores much larger amounts of energy: water is pumped to an upper reservoir
and later released through a reversible turbine. Its lifetime is long, but suitable terrain
and major civil works are required.

Compressed air, thermal storage and hydrogen address longer durations. Hydrogen is produced
by electrolysis, stored, and later used in a fuel cell or turbine. Because each conversion
loses energy, the round-trip efficiency is lower than that of a battery. Nevertheless,
hydrogen may become useful when storage duration matters more than efficiency, particularly
for seasonal balancing or industrial fuel supply.

No storage option is optimal for every task. A resilient low-carbon grid therefore combines
storage with interconnections, demand response, controllable generation and accurate
forecasting.
\end{ecarticle}

\begin{tcolorbox}[title=\sffamily\bfseries Did You Know?,colback=yellow!10,colframe=ECOrange]
A storage system can have excellent efficiency but still be unsuitable if its discharge
duration, response time, lifetime or location does not match the service required by the grid.
\end{tcolorbox}

\columnbreak

\begin{engineeringfocus}
\textbf{Pumped-storage hydroelectricity}

\begin{center}
\begin{tikzpicture}[scale=.72,transform shape,>=Latex,font=\scriptsize]
\fill[ECBlue!18] (0,2.7) rectangle (2.4,3.35);
\fill[ECBlue!18] (5.2,.1) rectangle (7.8,.75);
\draw[thick] (0,2.7) rectangle (2.4,3.35) node[midway]{upper reservoir};
\draw[thick] (5.2,.1) rectangle (7.8,.75) node[midway]{lower reservoir};
\node[draw=ECBlue,fill=ECLightBlue,rounded corners,minimum width=20mm,minimum height=9mm] (m) at (3.8,1.65) {motor--generator};
\node[draw=ECGreen,fill=ECGreen!10,rounded corners,minimum width=17mm,minimum height=9mm] (t) at (3.8,.65) {pump--turbine};
\draw[very thick,ECBlue] (2.4,3.0) -- (3.8,2.0) -- (t.north);
\draw[very thick,ECBlue] (t.east) -- (5.2,.45);
\draw[<->,very thick,ECOrange] (m) -- (t);
\draw[->,very thick,ECGreen] (2.8,2.35) -- node[above,sloped]{generation} (4.7,1.0);
\draw[->,very thick,ECPurple] (4.8,.82) -- node[below,sloped]{pumping} (2.7,2.2);
\end{tikzpicture}
\end{center}

During low demand, electrical energy drives the pump and raises water. During high demand,
the water flows downhill and drives the turbine. The stored gravitational energy is
$E=mgh$.
\end{engineeringfocus}

\begin{keyfigures}
\begin{tabularx}{\linewidth}{@{}Xr@{}}
Lithium-ion round-trip efficiency & 85--95\%\\
Pumped hydro & 70--85\%\\
Hydrogen power-to-power & 30--45\%\\
Typical battery duration & 1--8 h\\
Seasonal storage & weeks--months
\end{tabularx}
\end{keyfigures}

\begin{tcolorbox}[title=\sffamily\bfseries Storage technologies compared,colback=white,colframe=ECBlue]
\begin{center}
\begin{tikzpicture}
\begin{axis}[
width=\linewidth,height=48mm,
xbar,bar width=5pt,
xmin=0,xmax=100,
xlabel={Representative round-trip efficiency (\%)},
symbolic y coords={Hydrogen,Compressed air,Pumped hydro,Flywheel,Lithium-ion},
ytick=data,
yticklabel style={font=\scriptsize},
xticklabel style={font=\scriptsize},
grid=major]
\addplot[fill=ECBlue] table[x=efficiency,y=technology,col sep=comma]{data/EC05/storage_comparison.csv};
\end{axis}
\end{tikzpicture}
\end{center}
\footnotesize Representative values only: actual performance depends on design and operating conditions.
\end{tcolorbox}

\begin{vocabulary}
\begin{tabularx}{\linewidth}{@{}lX@{}}
round-trip efficiency & rendement aller-retour\\
charge / discharge & charge / décharge\\
storage duration & durée de stockage\\
reservoir & réservoir\\
flywheel & volant d'inertie\\
demand response & effacement de consommation\\
grid frequency & fréquence du réseau
\end{tabularx}
\end{vocabulary}

\begin{questions}
\item Why must electricity production and consumption remain balanced?
\item Distinguish power from stored energy.
\item Why are batteries suitable for rapid grid services?
\item Explain the operating cycle of a pumped-storage plant.
\item Why can hydrogen be useful despite its lower efficiency?
\item Which criteria should engineers consider besides efficiency?
\end{questions}

\begin{engineeringchallenge}
A remote island has a 12~MW peak demand, a 20~MW wind farm and a 10~MW solar farm.
Propose a storage strategy for frequency control, night-time supply and a five-day period
of weak wind. Select at least two technologies and justify their power, capacity, duration,
efficiency, safety and environmental impact.
\end{engineeringchallenge}

\begin{applicationexercise}
A pumped-storage plant transfers $5.0\times10^6$~m$^3$ of water through a vertical height of
420~m. Estimate the stored gravitational energy using $E=\rho Vgh$, with
$\rho=1000~\mathrm{kg\,m^{-3}}$. Express the result in GWh, then estimate the recoverable
energy for a round-trip efficiency of 78\%.
\end{applicationexercise}

\begin{speaking}
In teams, defend one storage technology in a three-minute technical pitch. Then identify
one service for which your technology would be a poor choice.
\end{speaking}
\end{multicols}

\subsection*{Sources}
\ECsource{IEA -- Grid-scale storage}{https://www.iea.org/energy-system/electricity/grid-scale-storage}
\ECsource{NREL -- Energy storage}{https://www.nrel.gov/grid/energy-storage.html}
